\lesson{2}{Sep 17 2021 Fri (11:40:03)}{What is Citizenship}{Unit 1}

\subsubsection*{Who is a Citizen}

The first part of the 14th Amendment to the United States Constitution defines
Citizens as people either \bf{"born or naturalized in the United States"}.
It then describes that Citizens are due certain protections and privileges
under the law that no state government may take away.

A person can be a United States Citizen at the time of birth in two main ways.
The Latin Term \bf{Jus Soli} means \bf{"Law of the Soil"}. People born
in the United States or United States territories are citizens because they
were born on United States Soil. An exception to the \bf{Jus Soli} rule
would be the child of a foreign diplomat as they are not considered subject to
the United States jurisdiction, a requirement of \bf{Jus Soli}.

The Latin Term \bf{Jus Sanguinis} means \bf{"law of the blood"}.

People born to at least one parent who is a United States Citizen and has lived
in the United States are citizens through their blood relationship to this
parent. United States Military and diplomatic sites in foreign countries are
not considered United States soil, contrary to what many believe. The
\bf{Jus Sanguinis} rule covers children born to United States Citizens
living abroad, such as military families. Most United States Citizen births,
but not all, meet both the \bf{Jus Soli Sanguinis} rules. The United States
has laws to define special situations.

\begin{figure}[h]
    \begin{quote}
        \it{All persons born or naturalized in the United States, and subject to
        the jurisdiction thereof, are citizens of the United States and of the State
        wherein they reside. No State shall make or enforce any law which shall abridge
        the privileges or immunities of citizens of the United States; nor shall any
        State deprive any person of life, liberty, or property, without due process of
        law; nor deny to any person within its jurisdiction the equal protection of the
        laws.}
    \end{quote}
    \caption{$14$th Amendment}
\end{figure}



\subsubsection*{How Does the Government Categorize Non-Citizens}

People living legally in the United States who are not United States Citizens
could have a visa or permanent resident status. They must carry proof of their
status at all times. A United States National is a person under the legal
protection of the United States but without citizenship. This applies to those
born United States territories like American Samoa and Swains Island. United
States Nationals can live and work in the United States but, at present, may
not vote in state or federal elections. They still have to go through full
naturalization process if they wish to become United States Citizens. However,
federal laws now grant full citizenship to people born in the territories of
Puerto Rico, the United States Virgin Island, and Guam.

A vis a gives a person permission to travel to a country and request entry by the country's officials. United States visas come in two basic types:

\begin{itemize}
  \item Non-immigrant visias
  \item Immigrant visas
\end{itemize}

People that wish to visit, go to school, or work temporarily in the United
States would apply for non-immigrant visa. This type of visa has a time limit.
Immigrant visas are for those who wish to live permanently in the United States
with no time limit. A person could apply for the immigrant visa and permanent
resident status at the same time.

Permanent resident status, formerly called resident alien, could take months or
years to obtain. This status gives the holder permission to live and work in
the United States. The United States permanent residency card is the official
document proving this permission. People often call it a \bf{\it{"green
card"}} because of the green ink used in the card's design. A permanent
resident must always carry the green card and pay income taxes to the United
States government. The United States still classifies permanent residents as
resident aliens. They do not have the full benefits of citizenship. However,
this status is an important prerequisite to applying for naturalization. A
person can apply for naturalization after being a permanent resident for a
certain length of time, usually five years.

\subsubsection*{What Are the Requirements of Naturalization}

You must do 4 things:

\begin{itemize}
  \item Application
  \item Fingerprints
  \item Interview
  \item Oath of Loyalty
\end{itemize}

\newpage
